% !PRINTED 4
\weberpage{4}

%% [running head, not body text -- commented out 2026-08-07] Page furniture (running
%% head, folio, signature) is excluded from this edition by standing policy.
%% \noindent{\small 4 \hfill Erster Abschnitt. \hfill \S.\,1.}

%<<CONT
und hieraus folgt durch die Schlussweise der vollst\"andigen Induction, dass man immer zu demselben Resultate kommt, wenn man in einer beliebigen Reihe von Elementen aus $P$ in endlicher Anzahl, $a, b, c, d \ldots$ zuerst zwei benachbarte Elemente componirt, dann wieder zwei benachbarte u.\,s.\,w., bis die ganze Reihe auf ein Element reducirt ist, das mit $abcd \ldots$ bezeichnet wird. So ist z.\,B.:
\[
\begin{aligned}
abcd &= (ab)cd = [(ab)c]d = (ab)(cd)\\
&= a(bc)d = [a(bc)]d = a[(bc)d]\\
&= ab(cd) = (ab)(cd) = a[b(c.d)]^{1)}. % [sic: the source prints a dot between c and d]
\end{aligned}
\]

\noindent 3.\ \gesperrt{Es wird vorausgesetzt. dass, wenn $ab = ab'$ oder $ab = a'b$ ist, nothwendig im ersten Falle $b = b'$, im zweiten $a = a'$ sein muss.} % [sic: the source prints a full stop after vorausgesetzt]

Wenn $P$ eine endliche Anzahl von Elementen umfasst, so heisst die Gruppe eine \gesperrt{endliche} und die Anzahl ihrer Elemente ihr \gesperrt{Grad}.

F\"ur endliche Gruppen ergiebt sich aus 1., 2., 3. die \gesperrt{Folgerung:}

\noindent 4.\ \gesperrt{Wenn von den drei Elementen $a, b, c$ aus $P$ zwei beliebig gegeben sind, so kann man das dritte immer und nur auf eine Weise so bestimmen, dass}
\[
ab = c
\]
\gesperrt{ist.}

Sind n\"amlich $a, b$ die gegebenen Elemente, so f\"allt die Behauptung 4. mit 1. zusammen. Ist aber $a$ und $c$ gegeben, so lasse man in dem Compositum $ab$ die zweite Componente $b$ das ganze System $P$ durchlaufen, dessen Grad $= n$ sei. Dann erh\"alt man nach 1. und 3. in $ab$ lauter verschiedene Elemente von $P$, und da ihre Anzahl $= n$ ist, so m\"ussen alle Elemente von $P$, also auch $c$, darunter vorkommen. Ebenso schliesst man, wenn $b$ und $c$ gegeben sind, indem man $a$ das ganze System $P$ durchlaufen l\"asst.

F\"ur unendliche Gruppen kann nicht mehr so geschlossen werden. \gesperrt{F\"ur unendliche Gruppen wird also noch die Eigenschaft 4. als Forderung in die Begriffsbestimmung mit aufgenommen.}

\noindent\rule{3cm}{0.4pt}

\noindent{\footnotesize $^{1)}$ Vgl.\ \gesperrt{Dirichlet-Dedekind}, Vorlesungen \"uber Zahlentheorie, \S.\,2.}
